% Part: lambda-calculus
% Chapter: lambda-definability
% Section: fixpoints

\documentclass[../../../include/open-logic-section]{subfiles}

\begin{document}

\olfileid{lam}{ldf}{fp}

\olsection{ثابت ټکي}

فرض کړئ چې د فکتوريال تابع د داسې ترم $\fn{Fac}$ په توګه
په بازګشتي ډول تعريفول غواړو چې دا خاصيت ولري:
\[
\fn{Fac} \ident \lambd[n][\fn{IsZero}\, n\, \num 1 (\fn{Mult}\, n (\fn{Fac}(\fn{Pred} \, n)))]
\]
يعنې که $n = 0$ وي، د $n$ فکتوريال $1$ دے؛ او که نۀ وي،
د $n$ او د $n-1$ د فکتوريال حاصل‌ضرب دے۔ خو ترم
$\fn{Fac}$ په دې ډول نۀ شو تعريفولے، ځکه $\fn{Fac}$
پخپله د تعريف ښي لوري ته راځي۔ د ځان‌مراجعې دغه ډول
بازګشتي تعريفونه د لامبډا حساب د ترمونو له عادي
تعريفونو څخه نۀ دي۔ د بېلګې په توګه، د يو ترم تعريف
داسې کېدے شي:
\[
\fn{Mult} \ident \lambd[ab][a (\fn{Add}\, b) 0]
\]
په ښي لوري کښې يوازې مخکې تعريف شوي ترمونه، لکه
$\fn{Add}$، کاروي۔ $\fn{Add}$ تل د خپل تعريفوونکي
ترم په ايښودلو لرې کولے شو۔ په دې توګه، د عددنښې
لڼدې ليکنې له پرانيستلو سره، $\fn{Mult}$ د خالص
لامبډا ترم بڼه خپلوي۔ که $\fn{Add}$ هم نور تعريف شوي
ترمونه کاروي، لکه $\fn{Add}'$، دغه پرله‌پسې ځاے پرځاے
کول غځوو او په پاې کښې خالص لامبډا ترم ته رسېږو۔
\textbf{د سرچينې سمون (OLLAM-055):} د ضرب د دې
بېلګې په سرچينه کښې دويم آرګومېنټ نۀ کارېږي؛ لکه
د مخکنۍ برخې په هم‌مانا فورمول کښې، دلته تکرارېدونکے
جمع کېدونکے دويم آرګومېنټ شوے دے۔

خو د پورته $\fn{Fac}$ په څېر بازګشتي تعريفونو کښې دا
چاره نۀ ترسره کېږي۔ که په ښي لوري کښې د $\fn{Fac}$
پېښه پخپله د $\fn{Fac}$ په تعريف بدله کړو، لاندې
عبارت ترلاسه کېږي:
\begin{align*}
  \fn{Fac} & \ident \lambd[n][\fn{IsZero}\, n\, \num{1}] \\
    & \qquad (\fn{Mult}\, n
  ((\lambd[n][\fn{IsZero} \, n \, \num 1\, (\fn{Mult}\, n\, (\fn{Fac}
    (\fn{Pred}\, n)))]) (\fn{Pred}\, n)))
\end{align*}
\textbf{د سرچينې سپيناوی (OLLAM-058):} په پورته
ښودنيزه پرانيستنه کښې د لومړۍ کرښې تړون له
بازګشتي څانګې څخه بېل چاپ شوے؛ د دواړو کرښو
ګډ مقصد د فکتوريال د ترم يوه پېښه پر خپل تعريف بدلول
دي، خو په دويمه کرښه کښې هم د ځان‌مراجعې نښه پاتې وي۔
او $\fn{Fac}$ لا هم په ښي لوري کښې پاتې دے۔ که همدا
کار بيا بيا وکړو، عبارت اوږدېږي او هېڅکله خالص
لامبډا ترم نۀ راکوي۔ نو د فکتوريال، يا په عام ډول د
بازګشتي تابعو، تعريف په دې لاره ممکن نۀ دے۔

خو بازګشتي تعريف يوه ګټوره خبره راښيي: که $f$ د
فکتوريال تابع تمثيلوونکے ترم وي، نو لاندې ترم
\[
\fn{Fac}' \ident \lambd[g][\lambd[n][\fn{IsZero} \, n \, \num 1\, (\fn{Mult} \, n\, (g (\fn{Pred} n)))]]
\]
چې پر $f$ تطبيق شي، يعنې $\fn{Fac}'\,f$، هم د
فکتوريال تابع تمثيل کوي۔ که $\fn{Fac}'$ د يوې تابعې
اخيستونکې او بله تابع ورکوونکې عمليې په توګه وګورو،
د فکتوريال پر يوه تمثيل ئې پايله بيا د هماغې تابعې
تمثيل وي۔ داسې ترم~$f$ چې $\fn{Fac}' \, f \equal[\beta] f$
خاصيت ولري، د $\fn{Fac}'$ \emph{ثابت ټکے} بلل کېږي۔
سرچينه وايي چې د $\fn{Fac'}\, f$ ارزښت هماغه $f$
دے، که $f$ د فکتوريال دپاره وي؛ په دې کښې د ترمونو
بېټا-برابري جلا ثبوت غواړي۔ سرچينه د فکتوريال تابع
په همدې معنا د $\fn{Fac}'$
ثابت ټکے بولي۔ \textbf{د سرچينې سپيناوی (OLLAM-056):}
د دوو ترمونو له دې خبرې چې دواړه يوه عددي تابع
محاسبه کوي، د هغو بېټا-برابري پخپله نۀ لازمېږي؛
د ثابت ټکي معادله يو جلا شرط دے۔

په لامبډا حساب کښې داسې ترمونه شته چې د ورکړل شوي
ترم ثابت ټکي پيدا کوي۔ په هغو سره د $\fn{Fac}'$ په
څېر ترم د فکتوريال په تعريف بدلولے شو۔

\begin{defn}
  \ollabel{defn:Turing-Y}
  \emph{Y-ترکيب کوونکے} دا ترم دے:
  \[
  Y \ident (\lambd[ux][x(uux)])(\lambd[ux][x(uux)]).
  \]
\end{defn}

\begin{thm}
  د هر ترم $g$ دپاره $Y$ دا خاصيت لري چې $Yg \red g(Yg)$۔
  نو $Yg$ تل د~$g$ ثابت ټکے دے۔
\end{thm}

\begin{proof}
  $(\lambd[ux][x(uux)])$ په~$U$ لنډوو، نو $Y \ident UU$۔ بيا
  \begin{align*}
    Y g &\ident (\lambd[ux][x(uux)])U\,g \\
    &\red (\lambd[x][x(UUx)])g\\
    & \red g(UUg) \ident g(Yg).
  \end{align*}
  څرنګه چې $g(Yg)$ او $Yg$ دواړه $g(Yg)$ ته
  راکمېږي، نو $g(Yg) \equal[\beta] Yg$؛ له دې امله
  $Yg$ د~$g$ ثابت ټکے دے۔
\end{proof}

څرنګه چې $Yg$ راکمېدونکے عبارت دے، راکمول ئې
بې‌پاېه دوام موندلے شي:
\begin{align*}
  Y g &\red g (Y g) \\
    &\red g (g (Y g)) \\
    &\red g(g (g (Y g)))\\
    &\ldots
\end{align*}
نو $Yg$ داسې انګېرلے شو لکه $g$ چې بې‌شمېره ځله
پر خپل تېر حاصل تطبيقېږي۔ که $g$ پر $Yg$ يو ځل
نور هم تطبيق کړو، په دې شهودي انځور کښې ګواکې
څۀ نۀ زياتېږي: $g$ پر $g$ تطبيقېږي، خو بيا هم
د $Yg$ د بې‌شمېره تکرار هماغه لړۍ ده؛ د $g$ يوه
نوره کارونه پکې ورکېږي۔ دا يوازې شهودي بيان دے؛
پورته کره دعوه د راکمېدنې او بېټا-برابرۍ ده۔

پام وکړئ چې له~$Yg$ څخه پيلېدونکې د $\beta$-راکمول
پورته لړۍ بې‌پاېه ده۔ که $Yg$ پر کوم ترم تطبيق کړو،
يعنې $(Yg)N$ وګورو، دا هم بې‌شمېره بېلابېلو ترمونو
ته راکمېدے شي: $(g(Yg))N$، $(g(g(Yg)))N$، \dots۔
سره له دې، کېدے شي د راکمولو کومه \emph{بله} لړۍ
په عادي بڼه پاے ته ورسېږي۔

د فکتوريال بېلګه واخلئ۔ $\fn{Fac}$ د $Y\, \fn{Fac}'$
په توګه تعريف کړئ، يعنې د $\fn{Fac'}$ يو ثابت ټکے۔ بيا:
\begin{align*}
  \fn{Fac}\,\num 3
  &\red Y\, \fn{Fac}' \, \num 3 \\
  &\red \fn{Fac}' (Y \,\fn{Fac}') \, \num 3 \\
  &\ident (\lambd[x][\lambd[n][\fn{IsZero} \, n\, \num 1\,
      (\fn{Mult}\, n\, (x (\fn{Pred}\, n)))]]) \, \fn{Fac} \, \num 3\\
  & \red \fn{IsZero} \, \num 3\, \num 1\,
      (\fn{Mult}\, \num 3\, (\fn{Fac} (\fn{Pred}\, \num 3))) \\
  &\red  \fn{Mult}\, \num 3 \, (\fn{Fac} \, \num 2).
  \intertext{په همدې ډول،}
  \fn{Fac}\,\num 2
  &\red  \fn{Mult}\, \num 2 \, (\fn{Fac} \, \num 1) \\
  \fn{Fac}\,\num 1
  &\red  \fn{Mult}\, \num 1 \, (\fn{Fac} \, \num 0)
  \intertext{خو}
  \fn{Fac}\,\num 0
  &\red \fn{Fac}' (Y \,\fn{Fac}') \, \num 0 \\
  &\ident (\lambd[x][\lambd[n][\fn{IsZero} \, n\, \num 1\,
      (\fn{Mult}\, n\, (x (\fn{Pred}\, n)))]]) \, \fn{Fac} \, \num 0\\
  & \red \fn{IsZero} \, \num 0\, \num 1\,
      (\fn{Mult}\, \num 0\, (\fn{Fac} (\fn{Pred}\, \num 0))).\\
  &\red \num 1.
  \intertext{نو په ټوليز ډول}
  \fn{Fac}\,\num 3
  &\red  \fn{Mult}\, \num 3 \,
  (\fn{Mult} \, \num 2\, (\fn{Mult} \, \num 1\, \num 1)).
\end{align*}

د $\fn{Fac'}$ په اړه دا چلند د هر بازګشتي تعريف
دپاره هم کار کوي۔ فرض کړئ لاندې بازګشتي معادله لرو:
\begin{align*}
g\,x_1\dots x_n & \equal[\beta] N
\intertext{چې $N$ ښايي~$g$ او $x_1$، \dots،~$x_n$ ولري۔ بيا تل داسې
ترم~$G \ident (Y \lambd[g][\lambd[x_1\dots x_n][N]])$ شته چې}
G\,x_1 \dots x_n & \equal[\beta] \Subst{N}{G}{g}.
\intertext{ځکه د ثابت ټکي د قضيې له مخې،}
G \ident (Y \lambd[g][\lambd[x_1\dots x_n][N]]) & \red \lambd[g][\lambd[x_1\dots x_n][N]](Y \lambd[g][\lambd[x_1\dots x_n][N]])\\
& \ident (\lambd[g][\lambd[x_1\dots x_n][N]])\,G
\intertext{او په پايله کښې}
G\,x_1 \dots x_n 
& \red (\lambd[g][\lambd[x_1\dots x_n][N]])\,G\,x_1\dots x_n\\
& \red (\lambd[x_1\dots x_n][\Subst{N}{G}{g}])\,x_1\dots x_n\\
  & \red \Subst{N}{G}{g}.
\end{align*}

د \olref{defn:Turing-Y} د $Y$ ترکيب کوونکے د الن
تورينګ دے۔ الونزو چرچ بله بڼه وړاندې کړې وه؛
هغې ته~$Y_C$ وايو:
\[
Y_C \ident \lambd[g][(\lambd[x][g(xx)])(\lambd[x][g(xx)])].
\]
د چرچ ترکيب کوونکے په يوه معنا د تورينګ له بڼې
څخه کمزورے دے: $Y_Cg \equal[\beta] g(Y_Cg)$، خو
$Y_Cg \bred g(Y_Cg)$ نۀ دے۔ \textbf{د سرچينې
سمون (OLLAM-057):} د دې دوو پرتلو په سرچينه کښې
د چرچ بڼې کښته نښه غورځېدلې وه؛ هغه دلته په دواړو
ځايونو کښې څرګنده شوې ده۔ $V$ دې د
$\lambd[x][g(xx)]$ ترم وي، نو $Y_C \ident \lambd[g][VV]$۔ بيا
\begin{align*}
  VV & \ident (\lambd[x][g(xx)])V \red g(VV)
  \text{ او له دې امله}\\
  Y_C g & \ident (\lambd[g][VV])g \red VV \red g(VV), \text{ خو همدارنګه}\\
  g(Y_C g) & \ident g((\lambd[g][VV])g) \red g(VV).
\end{align*}
په بله وينا، $Y_Cg$ او $g(Y_Cg)$ دواړه ګډ ترم
$g(VV)$ ته راکمېږي؛ نو $Y_Cg \equal[\beta] g(Y_Cg)$۔
د ډېرو کارونو دپاره همدا بس وي۔

\end{document}
